\section{模型评价与推广}

\subsection{模型的适用范围与优势}

本文四问共用一套柱坐标径向热-质耦合方程，问题之间的差别只落在物性闭合式、求解域与终止判据三处。
这种组织方式带来两个直接好处：其一，问题一的半解析对照一旦通过，其验证的是同一套有限体积离散核，
因此该结论可以顺延到后三问的温度场；其二，问题四补齐 $2\times2$ 因子后可以同时看到物性路径、
收缩路径及其交互，而不是把 $-5.926$~h 的净差说成单一原因。

分布参数描述的选择也有定量依据。表面对流与内部导热、内部扩散之比 $Bi_h=1.389$、$Bi_m=3.240$
均为 $O(1)$ 量级，既排除了把药材当作等温等湿集中参数体的简化，也排除了把表面直接固定为环境值的
第一类边界。灵敏度探针给出的对照支持这一判断：改用第一类边界后烘干时长由 $56.18$~h 缩短至
$51.23$~h，偏差约 $8.8\%$，远超其他数值设置带来的差异。移动边界采用物质坐标而非在物理坐标下
重构网格，使收缩过程中网格拓扑保持不变；$\dot R\to 0$ 时退化为固定域，与 Eulerian 固定域
时长相对差 $0.24\%$，来自结点分布而不是额外物理项。

\subsection{模型特点}

第一，问题四用物质坐标把固相速度写清楚，而不是把对流项当作可开可关的数值开关。
均匀收缩给出 $v_s=r\dot R/R$，干基含水率跟随物质点，交付时长为 $50.780$~h；
实验室参考系下的 Landau 表述给出 $52.190$~h，相对差 $2.78\%$，作为坐标表述的模型区间一并报告。

第二，达标判据按题目「各处」取逐点最大值。以体平均判定时，$\bar C$ 会在芯部远未达标时先跨过阈值：
$C_{\max}$ 恰触 $0.15$~kg/kg 的时刻体平均仅 $0.1232$~kg/kg。判据的选择在这里改变的是答案本身。

第三，物性与收缩的贡献用完整 $2\times2$ 报告，并给出交互项 $I=-46.240$~h。
两条路径的 Shapley 平均为物性 $+48.812$~h、收缩 $-55.006$~h，但交互项与主效应同量级，
因此不把加性分解说成唯一归因。

第四，附录4 的 $\rho(C)$ 与附件2 的 $R(t)$ 在严格干物质守恒下相差 $13.42\%$。
本文依题面把实测半径作为给定边界，把不相容度降格为诊断量，并且不用坐标交叉验证的
$2.78\%$ 去代替该项对时长的影响。

\subsection{模型的局限}

第一项局限来自能量方程中未设蒸发潜热汇。药材内部水分汽化要吸收汽化潜热，忽略这一项会高估内部
温度，进而高估水分扩散系数，使预测的烘干时长偏短。§\ref{sec:sens} 的结构性探查量化了这一偏向：
在能量方程中按 $2.4\times 10^{6}$~J/kg 计入潜热汇后，基准时长由 $56.18$~h 延长至 $62.12$~h，
增幅约 $10.6\%$。因此本文给出的 $56.706$~h 与 $50.780$~h 应当理解为乐观的短侧估计，工艺上若以此
安排烘干批次，需要预留一成左右的时间余量。

第二项局限是均匀收缩闭合。式~\eqref{eq:vs} 假定骨架按仿射方式运动，题目没有内部位移观测来
证实或否定这一点。Eulerian 交叉验证把由此带来的时长差限制在 $2.78\%$ 量级，但不能排除
非均匀塌缩。

第三项局限是附录4 密度式与附件2 实测半径不相容。由密度反推的末态半径为 $1.287$~cm，
实测为 $1.198$~cm。交付结果使用实测 $R(t)$，该数据冲突没有被「修正掉」，其对时长的影响
也没有被单独重算。

第四项局限在于灵敏度分析采用单因素扰动，且验证以内部自洽为主。半解析解、网格收敛和离散收支
残差证明的是数值实现正确，不能写成「充分验证了模型真实正确」。缺少药材内部温湿度实测剖面，
模型参数未经标定。

\subsection{结论与推广}

按附录2 的常物性闭合，$1800$~s 预热段结束时中心温度升至 $33.577$~℃ 而中心含水率仍为
$2.5500$~kg/kg，表面已降至 $1.511$~kg/kg，热扩散与质扩散的时间尺度之比约 $34$，“热先于质”由此得到
定量表述。按附录3 的变物性强耦合闭合，$3$~h 时中心温度达 $49.85$~℃ 已接近烘房温度，而中心含水率
$1.766$~kg/kg 距达标仍远，烘干时长为 $56.706$~h。计入附录4 物性与附件2 收缩规律后时长为
$50.780$~h，比问题三缩短 $5.926$~h；同物性下相对尺寸冻结缩短 $60.6\%$。
$2\times2$ 因子的交互项为 $-46.240$~h，物性与收缩不能作唯一加性分解。
仅凭“体积缩小所以干得快”或“扩散系数变小所以干得慢”任一单侧推理都会得到错误的量级判断。

本文的建模路线可以移用到其他伴随收缩的多孔介质干燥问题。只要给出与含水率相关的物性闭合以及可观测
的外形收缩数据，把物性式与 $R(t)$ 替换即可沿用同一求解核；若干燥对象为球形或平板，仅需把控制体
的面积权重由 $2\pi r$ 改为 $4\pi r^2$ 或常数，球形移动边界下的本征展开解可用作同类
对照\upcite{Ali2026}。若收缩规律本身未知而只有出口含水率可测，则问题转为由观测反求边界运动的
反问题，需要另一类求解策略\upcite{Lotfi2019}。对工艺参数的优化也可以在此框架内进行：由于
$D$ 前因子的弹性接近 $-1$ 而 $h$ 的弹性近于零，缩短烘干时间的有效手段在于提高介质温度以抬升扩散
系数，而不是加大风速强化表面换热。若进一步需要预测能耗或品质指标，需在能量方程中补入潜热汇并引入
温度-时间累积的品质衰减描述，这两项都超出题给信息，属于后续工作。

% DATA_CHECK_PASSED
